% ============================================================
% CAPITOLO 31 - LA REGOLA DELL'OTTAVA COMPLETA (DIATONICA)
% ============================================================
\chapter{La Regola dell'Ottava completa (diatonica)}
\label{chap:regola_ottava_completa}

Dopo la teoria delle settime si può presentare la \textbf{versione completa} della Regola dell'Ottava, che integra le settime diatoniche nell'armonizzazione della scala al basso. Qui si trattano solo accordi che stanno nei suoni della tonalità; le alterazioni cromatiche sono nel Capitolo \ref{chap:regola_ottava_cromatica}.

\section{Le tre versioni a confronto}
\label{sec:tre_versioni_confronto}

\begin{definizione}[Le tre versioni della Regola dell'Ottava]
	\label{def:tre_versioni_complete}
	\begin{enumerate}
		\item \textbf{Versione preliminare (Capitolo \ref{chap:ottava_preliminare}):} solo \textbf{triadi}, in stato fondamentale e rivoltate. L'armonizzazione è tutta consonante e serve a imparare la condotta delle voci.
		\item \textbf{Versione intermedia (Capitolo \ref{chap:settima_dominante}):} aggiunge la \textbf{settima di dominante (\sigla{V}{7}{})} e i suoi tre rivolti, una dissonanza controllata che dà al percorso più tensione e direzione.
		\item \textbf{Versione completa (questo capitolo):} include le \textbf{settime diatoniche}, cioè la \sigla{V}{7}{} con i suoi rivolti e il \sigla{II}{6}{5}. Offre più scelte di colore armonico e corrisponde alla pratica del tardo Barocco e del primo Classicismo.
	\end{enumerate}
\end{definizione}

Quanto cresce il numero delle scelte da una versione all'altra è calcolato nell'Approfondimento \ref{math:crescita_combinatoria}.

\section{La Regola dell'Ottava nel modo maggiore}
\label{sec:regola_completa_maggiore}

\begin{regola}[Regola dell'Ottava completa: modo maggiore]
	\label{reg:regola_ottava_completa_maggiore}

	\textbf{Scala ascendente (esempio in Do maggiore):}

	\begin{tabular}{cccl}
		\toprule
		\textbf{Grado} & \textbf{Basso} & \textbf{Numerica} & \textbf{Funzione armonica} \\
		\midrule
		I & Do & 5 & I (tonica) \\
		II & Re & 6 \textit{oppure} 4/3 & \sigla{VII}{6}{} \textit{oppure} \sigla{V}{4}{3} \\
		III & Mi & 6 & \sigla{I}{6}{} (tonica) \\
		IV & Fa & 6 \textit{oppure} 6/5 & \sigla{II}{6}{} \textit{oppure} \sigla{II}{6}{5} (sopratonica con settima) \\
		V & Sol & 5 \textit{oppure} 7 & V \textit{oppure} \sigla{V}{7}{} (dominante) \\
		VI & La & 6 & \sigla{IV}{6}{} (sottodominante) \\
		VII & Si & 6 \textit{oppure} 6/5 & \sigla{V}{6}{} \textit{oppure} \sigla{V}{6}{5} (dominante con settima) \\
		VIII & Do & 5 & I (tonica) \\
		\bottomrule
	\end{tabular}

	\vspace{0.5cm}

	\textbf{Scala discendente (esempio in Do maggiore):}

	\begin{tabular}{cccl}
		\toprule
		\textbf{Grado} & \textbf{Basso} & \textbf{Numerica} & \textbf{Funzione armonica} \\
		\midrule
		VIII & Do & 5 & I (tonica) \\
		VII & Si & 6 & \sigla{V}{6}{} (dominante) \\
		VI & La & 6 & \sigla{IV}{6}{} (sottodominante) \\
		V & Sol & 5 \textit{oppure} 7 & V \textit{oppure} \sigla{V}{7}{} (dominante) \\
		IV & Fa & 2 & \sigla{V}{2}{} (dominante, terzo rivolto) \\
		III & Mi & 6 & \sigla{I}{6}{} (tonica) \\
		II & Re & 4/3 & \sigla{V}{4}{3} (dominante, secondo rivolto) \\
		I & Do & 5 & I (tonica) \\
		\bottomrule
	\end{tabular}
\end{regola}

\begin{esempio}[Regola dell'Ottava completa in Sol maggiore: scala discendente]
	\label{ex:regola_completa_disc_magg}
	L'esempio è in Sol maggiore, perché la scala di Sol (Sol$_3$--Sol$_2$) tiene il basso nella zona preferibile; lo schema è quello della tabella per Do maggiore. Il VII grado ha il \sigla{V}{6}{}, il IV il \sigla{V}{2}{}, il II il \sigla{V}{4}{3}: la settima di dominante colora quasi tutta la discesa, mentre il soprano sale in moto contrario al basso.
	\musica{es_cap31_01}
\end{esempio}

\section{La Regola dell'Ottava nel modo minore}
\label{sec:regola_completa_minore}

Nel modo minore la versione completa ha alcune particolarità, legate alla scelta tra scala minore armonica e melodica. In salita la scala armonica fa fare al basso la seconda eccedente Fa--Sol$\sharp$ (Osservazione \ref{oss:seconda_aumentata}); la scala melodica la evita alzando anche il VI grado (Sezione \ref{sec:regola_completa_melodica}).

\begin{regola}[Regola dell'Ottava completa: modo minore]
	\label{reg:regola_ottava_completa_minore}

	\textbf{Scala ascendente (esempio in La minore armonico):}

	\begin{tabular}{cccl}
		\toprule
		\textbf{Grado} & \textbf{Basso} & \textbf{Numerica} & \textbf{Funzione armonica} \\
		\midrule
		I & La & 5 & i (tonica) \\
		II & Si & 4/3 & \sigla{V}{4}{3} (dominante) \\
		III & Do & 6 & \sigla{i}{6}{} (tonica) \\
		IV & Re & 6 \textit{oppure} 6/5 & \sigla{ii°}{6}{} \textit{oppure} \sigla{iiø}{6}{5} (sopratonica con settima) \\
		V & Mi & 5 \textit{oppure} 7 & V \textit{oppure} \sigla{V}{7}{} (dominante) \\
		VI & Fa & 6 & \sigla{iv}{6}{} (sottodominante) \\
		VII & Sol$\sharp$ & 6/5 & \sigla{V}{6}{5} (dominante con settima) \\
		VIII & La & 5 & i (tonica) \\
		\bottomrule
	\end{tabular}

	\vspace{0.5cm}

	\textbf{Scala discendente (esempio in La minore):}

	\begin{tabular}{cccl}
		\toprule
		\textbf{Grado} & \textbf{Basso} & \textbf{Numerica} & \textbf{Funzione armonica} \\
		\midrule
		VIII & La & 5 & i (tonica) \\
		VII & Sol$\sharp$ & 6 & \sigla{V}{6}{} (dominante) \\
		VI & Fa & 6 & \sigla{iv}{6}{} (sottodominante) \\
		V & Mi & 5 \textit{oppure} 7 & V \textit{oppure} \sigla{V}{7}{} (dominante) \\
		IV & Re & 2 & \sigla{V}{2}{} (dominante, terzo rivolto) \\
		III & Do & 6 & \sigla{i}{6}{} (tonica) \\
		II & Si & 4/3 & \sigla{V}{4}{3} (dominante, secondo rivolto) \\
		I & La & 5 & i (tonica) \\
		\bottomrule
	\end{tabular}
\end{regola}

\begin{esempio}[Regola dell'Ottava completa in La minore: scala ascendente]
	\label{ex:regola_completa_asc_min}
	Il II grado ha il \sigla{V}{4}{3}, il IV il \sigla{iiø}{6}{5}, il VII il \sigla{V}{6}{5}. Fra il VI e il VII grado il basso sale per seconda eccedente (Fa--Sol$\sharp$): la versione melodica dell'Esempio \ref{ex:regola_completa_asc_melodica} la evita.
	\musica{es_cap14_02}
\end{esempio}

\begin{esempio}[Regola dell'Ottava completa in La minore: scala discendente]
	\label{ex:regola_completa_disc_min}
	\musica{es_cap31_02}
\end{esempio}

\section{La scala minore melodica ascendente}
\label{sec:regola_completa_melodica}

Salendo, il basso del modo minore segue di solito la scala \textbf{melodica}: il VI e il VII grado sono alzati (Fa$\sharp$ e Sol$\sharp$ in La minore) e il tratto dal V all'VIII grado coincide con quello della scala maggiore. Scendendo si torna ai gradi naturali o a quelli della scala armonica, come nelle tabelle precedenti.

\begin{regola}[Regola dell'Ottava completa: scala minore melodica ascendente]
	\label{reg:regola_ottava_completa_melodica}

	\textbf{Scala ascendente (esempio in La minore melodica):}

	\begin{tabular}{cccl}
		\toprule
		\textbf{Grado} & \textbf{Basso} & \textbf{Numerica} & \textbf{Funzione armonica} \\
		\midrule
		I & La & 5 & i (tonica) \\
		II & Si & 4/3 & \sigla{V}{4}{3} (dominante) \\
		III & Do & 6 & \sigla{i}{6}{} (tonica) \\
		IV & Re & 6 \textit{oppure} 6/5 & \sigla{ii}{6}{} \textit{oppure} \sigla{ii}{6}{5} (sopratonica, con Fa$\sharp$) \\
		V & Mi & 5 \textit{oppure} 7 & V \textit{oppure} \sigla{V}{7}{} (dominante) \\
		VI & Fa$\sharp$ & 6 & \sigla{IV}{6}{} (sottodominante maggiore) \\
		VII & Sol$\sharp$ & 6/5 & \sigla{V}{6}{5} (dominante con settima) \\
		VIII & La & 5 & i (tonica) \\
		\bottomrule
	\end{tabular}

	\vspace{0.5cm}

	Sul VI grado alzato la numerica resta 6 (sesta Re e terza La, entrambe naturali): l'accordo è la triade maggiore di IV (Re-Fa$\sharp$-La) in primo rivolto. Il Fa$\sharp$ del basso sale alla sensibile e non si raddoppia; si raddoppia il Re o il La. Il Re, tenuto nella stessa voce, diventa la settima del \sigla{V}{6}{5} seguente, che così è preparata e scende sul Do della tonica. Sul IV grado il Fa$\sharp$ dà la triade minore di II (Si-Re-Fa$\sharp$) e, con la settima La, il \sigla{ii}{6}{5}, la cui settima scende sul Sol$\sharp$ del V.
\end{regola}

\begin{esempio}[Regola dell'Ottava completa in La minore melodica: scala ascendente]
	\label{ex:regola_completa_asc_melodica}
	Il basso sale La--Si--Do--Re--Mi--Fa$\sharp$--Sol$\sharp$--La senza seconde eccedenti. Il IV grado ha il \sigla{ii}{6}{5}, il VI il \sigla{IV}{6}{} maggiore, il VII il \sigla{V}{6}{5} con la settima preparata dal soprano.
	\musica{es_cap31_03}
\end{esempio}

\section{Il \sigla{V}{6}{5} e il \sigla{II}{6}{5} nella Regola}
\label{sec:rivolti_regola}

Due rivolti di settima compaiono nella regola su gradi del basso che non sono il loro grado di origine: il \sigla{V}{6}{5} sul VII grado e il \sigla{II}{6}{5} sul IV grado.

\begin{regola}[\sigla{V}{6}{5} sul VII grado nella scala ascendente]
	\label{reg:vii65_regola}
	Nella scala ascendente del modo maggiore, quando il basso raggiunge il VII grado, si usa spesso il primo rivolto della settima di dominante (\sigla{V}{6}{5}): sul basso Si stanno Re, Fa e Sol, cioè la triade di VII con la fondamentale Sol in più. La sensibile al basso sale alla tonica e la settima (Fa), in una voce superiore, scende di grado sul Mi della tonica. Nella scala discendente il VII grado si armonizza invece con il \sigla{V}{6}{}, perché il basso prosegue verso il VI e la settima non troverebbe la sua risoluzione.
\end{regola}

\begin{esempio}[\sigla{V}{6}{5} nella Regola dell'Ottava in Do maggiore]
	\label{ex:vii65_regola_maggiore}
	Frammento con il \sigla{V}{6}{5} sul VII grado, che risolve sulla tonica.
	\musica{es_cap12_02}
\end{esempio}

\begin{regola}[\sigla{II}{6}{5} sul IV grado nella scala ascendente del modo minore]
	\label{reg:ii65_regola_minore}
	Nella scala ascendente del modo minore, come nel maggiore, il IV grado al basso può essere armonizzato con il primo rivolto della settima del II grado (\sigla{II}{6}{5}). La settima è preparata e risolve scendendo di grado nell'accordo successivo, di norma il V.
\end{regola}

\begin{esempio}[\sigla{II}{6}{5} nella Regola dell'Ottava in La minore]
	\label{ex:ii65_regola_minore}
	Frammento della scala ascendente in La minore (III, IV e V grado al basso).
	\musica{es_cap12_05}
\end{esempio}

\section{Usi pratici e stilistici; tabella di riferimento}
\label{sec:usi_pratici_completi}

Fra le possibilità della versione completa si sceglie secondo il contesto musicale.

\begin{regola}[Criteri di scelta tra le versioni]
	\label{reg:criteri_scelta_versioni}
	\begin{itemize}
		\item \textbf{Periodo storico:} primo Barocco, versione preliminare; tardo Barocco (Bach, Händel), intermedia o completa; Classicismo (Mozart, Haydn), completa con libertà di scelta.
		\item \textbf{Contesto espressivo:} passaggi leggeri e danzanti, preliminare; momenti drammatici e intensi, intermedia o completa; sequenze modulanti, completa (le settime facilitano la modulazione).
		\item \textbf{Velocità del basso:} basso veloce (semicrome, terzine), triadi semplici; basso moderato (crome, semiminime), intermedia; basso lento (minime, semibrevi), completa.
		\item \textbf{Tessitura:} contrappunto denso, preliminare (si evitano troppe dissonanze); accompagnamento omofono, completa.
	\end{itemize}
\end{regola}

Perché le settime «colorano» l'armonia è detto nell'Approfondimento \ref{phys:settime_colore}.

\begin{table}[h]
	\centering
	\caption{Regola dell'Ottava completa: tabella di riferimento (versione diatonica)}
	\label{tab:regola_ottava_completa_riepilogo}
	\small
	\vspace{0.2cm}
	\textbf{Scala ascendente}\\[0.2cm]
	\begin{tabular}{cllllll}
		\toprule
		\textbf{Grado} & \multicolumn{2}{c}{\textbf{Do maggiore}} & \multicolumn{2}{c}{\textbf{La minore armonica}} & \multicolumn{2}{c}{\textbf{La minore melodica}} \\
		\cmidrule(lr){2-3}\cmidrule(lr){4-5}\cmidrule(lr){6-7}
		& Basso & Accordo & Basso & Accordo & Basso & Accordo \\
		\midrule
		I & Do & \sigla{I}{}{} & La & \sigla{i}{}{} & La & \sigla{i}{}{} \\
		II & Re & \sigla{vii°}{6}{} \textit{o} \sigla{V}{4}{3} & Si & \sigla{V}{4}{3} & Si & \sigla{V}{4}{3} \\
		III & Mi & \sigla{I}{6}{} & Do & \sigla{i}{6}{} & Do & \sigla{i}{6}{} \\
		IV & Fa & \sigla{ii}{6}{} \textit{o} \sigla{ii}{6}{5} & Re & \sigla{ii°}{6}{} \textit{o} \sigla{iiø}{6}{5} & Re & \sigla{ii}{6}{} \textit{o} \sigla{ii}{6}{5} \\
		V & Sol & V \textit{o} \sigla{V}{7}{} & Mi & V \textit{o} \sigla{V}{7}{} & Mi & V \textit{o} \sigla{V}{7}{} \\
		VI & La & \sigla{IV}{6}{} & Fa & \sigla{iv}{6}{} & Fa$\sharp$ & \sigla{IV}{6}{} \\
		VII & Si & \sigla{V}{6}{} \textit{o} \sigla{V}{6}{5} & Sol$\sharp$ & \sigla{V}{6}{5} & Sol$\sharp$ & \sigla{V}{6}{5} \\
		VIII & Do & \sigla{I}{}{} & La & \sigla{i}{}{} & La & \sigla{i}{}{} \\
		\bottomrule
	\end{tabular}

	\vspace{0.4cm}
	\textbf{Scala discendente}\\[0.2cm]
	\begin{tabular}{cllll}
		\toprule
		\textbf{Grado} & \multicolumn{2}{c}{\textbf{Do maggiore}} & \multicolumn{2}{c}{\textbf{La minore armonica}} \\
		\cmidrule(lr){2-3}\cmidrule(lr){4-5}
		& Basso & Accordo & Basso & Accordo \\
		\midrule
		VIII & Do & \sigla{I}{}{} & La & \sigla{i}{}{} \\
		VII & Si & \sigla{V}{6}{} & Sol$\sharp$ & \sigla{V}{6}{} \\
		VI & La & \sigla{IV}{6}{} & Fa & \sigla{iv}{6}{} \\
		V & Sol & V \textit{o} \sigla{V}{7}{} & Mi & V \textit{o} \sigla{V}{7}{} \\
		IV & Fa & \sigla{V}{2}{} & Re & \sigla{V}{2}{} \\
		III & Mi & \sigla{I}{6}{} & Do & \sigla{i}{6}{} \\
		II & Re & \sigla{V}{4}{3} & Si & \sigla{V}{4}{3} \\
		I & Do & \sigla{I}{}{} & La & \sigla{i}{}{} \\
		\bottomrule
	\end{tabular}

	\vspace{0.2cm}
	\parbox{0.9\textwidth}{\footnotesize La numerica è scritta accanto al grado di ogni sigla; \textit{o} separa le alternative. In salita la scala melodica evita la seconda eccedente Fa--Sol$\sharp$ del basso (Sezione \ref{sec:regola_completa_melodica}).}
\end{table}
